% sections/statistical_foundations.tex — Statistical Foundations
\section{Statistical Foundations}
\label{sec:statistical_foundations}

This section formalizes the statistical machinery underlying \system{}'s
data generation.  We cover Benford-compliant amount sampling, copula-based
dependency modeling, mixture distributions, and temporal dynamics.

% ─────────────────────────────────────────────────────────────────────
\subsection{Benford's Law Compliance}
\label{sec:stat_benford}

Benford's law~\citep{benford1938law,nigrini2012benfords} states that the
leading digit $d \in \{1, \ldots, 9\}$ of naturally occurring numerical data
follows the distribution:

\begin{equation}
P(D_1 = d) = \log_{10}\!\left(1 + \frac{1}{d}\right).
\label{eq:benford_first}
\end{equation}

The generalization to the first two digits $d_1 d_2 \in \{10, \ldots, 99\}$
is:

\begin{equation}
P(D_{12} = d_1 d_2) = \log_{10}\!\left(1 + \frac{1}{d_1 \cdot 10 + d_2}\right).
\label{eq:benford_two}
\end{equation}

\system{} implements three Benford-aware samplers:

\paragraph{BenfordSampler.}
Uses CDF-based inverse-transform sampling over the first-digit probabilities.
Given a pre-computed CDF vector
$F = [0.30103, 0.47712, 0.60206, \ldots, 1.0]$,
a uniform draw $u \sim \mathcal{U}(0,1)$ is mapped to digit $d$ via binary
search: $d = \min\{k : F[k] \geq u\}$.  The magnitude (order of magnitude)
is sampled adaptively from a configured range, and trailing digits are drawn
uniformly.

\paragraph{EnhancedBenfordSampler.}
Extends first-digit sampling to the first-two-digit distribution
(\Cref{eq:benford_two}), providing finer-grained statistical fidelity.  This
is critical for passing second-digit tests commonly used in forensic
accounting~\citep{nigrini2012benfords}.

\paragraph{BenfordDeviationSampler.}
Generates amounts that \emph{intentionally violate} Benford's law for anomaly
injection.  Five deviation modes are supported:

\begin{itemize}
  \item \textbf{RoundNumberBias}: Over-represents amounts ending in multiples
    of powers of ten (e.g., \$1{,}000, \$5{,}000).
  \item \textbf{ThresholdClustering}: Concentrates amounts just below
    approval thresholds (e.g., \$4{,}999 for a \$5{,}000 limit).
  \item \textbf{UniformFirstDigit}: Replaces Benford with $\mathcal{U}\{1,9\}$.
  \item \textbf{DigitBias}: Over-weights specific digits (anti-Benford
    probabilities emphasize 5, 7, 9).
  \item \textbf{TrailingZeros}: Forces amounts to have zero trailing digits.
\end{itemize}

\begin{proposition}[MAD bound]
\label{prop:mad}
Let $\hat{p}_d$ denote the empirical first-digit frequency from $n$ samples
of the BenfordSampler.  The mean absolute deviation
$\MAD = \frac{1}{9}\sum_{d=1}^{9}|\hat{p}_d - p_d|$
satisfies $\E[\MAD] \leq \frac{1}{9}\sqrt{\frac{\pi}{2n}}$ for large $n$,
ensuring convergence to ``close conformity'' (MAD~$< 0.006$) as classified
by~\citet{nigrini2012benfords} for $n \gtrsim 7{,}000$.
\end{proposition}

% ─────────────────────────────────────────────────────────────────────
\subsection{Copula-Based Dependency Modeling}
\label{sec:stat_copula}

Real accounting data exhibits complex dependencies: transaction amounts
correlate with line-item counts, approval levels depend on amounts, and
payment timing relates to invoice size.  \system{} separates marginal
distributions from dependency structure using copulas~\citep{nelsen2006introduction}.

\begin{definition}[Copula]
A $d$-dimensional copula $C : [0,1]^d \to [0,1]$ is a joint CDF with
uniform marginals.  By Sklar's theorem, any joint distribution $H$ with
marginals $F_1, \ldots, F_d$ can be written as
$H(x_1, \ldots, x_d) = C(F_1(x_1), \ldots, F_d(x_d))$.
\end{definition}

\system{} implements five copula families, each capturing different tail
dependency patterns relevant to financial data:

\paragraph{1. Gaussian Copula.}
Parameterized by correlation matrix $\bm{\Sigma}$.  No tail dependence;
symmetric.  Generated via Cholesky decomposition $\bm{L}\bm{L}^\top = \bm{\Sigma}$
applied to independent standard normals: $\bm{z} = \bm{L}\bm{w}$,
$\bm{w} \sim \mathcal{N}(\bm{0}, \bm{I})$.  Kendall's tau:
$\tau = \frac{2}{\pi}\arcsin(\rho)$.

\paragraph{2. Clayton Copula.}
$C_\theta(u,v) = (u^{-\theta} + v^{-\theta} - 1)^{-1/\theta}$, $\theta > 0$.
Exhibits \emph{lower tail dependence}, suitable for modeling joint
distress scenarios (e.g., correlated payment defaults).  Kendall's tau:
$\tau = \theta / (\theta + 2)$.  Sampling uses the conditional method.

\paragraph{3. Gumbel Copula.}
$C_\theta(u,v) = \exp\!\big({-}[(-\ln u)^\theta + (-\ln v)^\theta]^{1/\theta}\big)$,
$\theta \geq 1$.
Exhibits \emph{upper tail dependence}, appropriate for correlated high-value
transactions.  Uses the Marshall-Olkin algorithm with a positive stable
distribution.  $\tau = 1 - 1/\theta$.

\paragraph{4. Frank Copula.}
$C_\theta(u,v) = -\frac{1}{\theta}\ln\!\left(1 + \frac{(e^{-\theta u}-1)(e^{-\theta v}-1)}{e^{-\theta}-1}\right)$.
No tail dependence; symmetric.  Sampling via conditional distribution inversion.
Kendall's tau involves the Debye function $D_1$.

\paragraph{5. Student-$t$ Copula.}
Parameterized by correlation $\rho$ and degrees of freedom $\nu$.  Exhibits
\emph{symmetric tail dependence}, capturing joint extreme events.  Generated
via $\bm{z} = \bm{L}\bm{w}/\sqrt{S/\nu}$ where $S \sim \chi^2_\nu$.

\begin{table}[t]
\centering
\caption{Copula properties and recommended accounting applications.}
\label{tab:copula_summary}
\small
\begin{tabular}{lccl}
\toprule
\textbf{Copula} & \textbf{Lower Tail} & \textbf{Upper Tail}
  & \textbf{Use Case} \\
\midrule
Gaussian    & None & None & General correlation \\
Clayton     & Yes  & None & Payment defaults, distress \\
Gumbel      & None & Yes  & High-value co-occurrence \\
Frank       & None & None & Symmetric, no extremes \\
Student-$t$ & Yes  & Yes  & Joint tail events \\
\bottomrule
\end{tabular}
\end{table}

% ─────────────────────────────────────────────────────────────────────
\subsection{Mixture Distributions for Transaction Amounts}
\label{sec:stat_mixture}

Enterprise transaction amounts naturally cluster into distinct regimes
(routine, significant, major).  \system{} models this via $K$-component
log-normal mixtures:

\begin{equation}
f(x) = \sum_{k=1}^{K} w_k \cdot \frac{1}{x\sigma_k\sqrt{2\pi}}
  \exp\!\left(-\frac{(\ln x - \mu_k)^2}{2\sigma_k^2}\right),
  \quad \sum_{k=1}^{K} w_k = 1.
\label{eq:lognormal_mixture}
\end{equation}

The default three-component configuration mimics empirical distributions
observed in ERP data:

\begin{center}
\begin{tabular}{lcccc}
\toprule
\textbf{Component} & $w_k$ & $\mu_k$ & $\sigma_k$
  & \textbf{Mean} ($e^{\mu+\sigma^2/2}$) \\
\midrule
Routine     & 0.60 & 6.0 & 1.5 & \$1{,}224 \\
Significant & 0.30 & 8.5 & 1.0 & \$8{,}103 \\
Major       & 0.10 & 11.0 & 0.8 & \$89{,}352 \\
\bottomrule
\end{tabular}
\end{center}

Component selection uses pre-computed cumulative weights with $O(\log K)$
binary search.  Gaussian mixtures are also supported for fields with
real-valued (possibly negative) domains such as foreign-exchange translation
adjustments.

% ─────────────────────────────────────────────────────────────────────
\subsection{Temporal Dynamics}
\label{sec:stat_temporal}

Financial data exhibits rich temporal structure that must be faithfully
reproduced.

\paragraph{Business-day calendars.}
\system{} implements holiday calendars for 11 regions (US, DE, GB, CN, JP,
IN, BR, MX, AU, SG, KR) with configurable half-day policies, month-end
conventions (modified following, preceding, end-of-month), and T+$N$
settlement rules.

\paragraph{Period-end dynamics.}
Transaction volumes spike near period ends.  We model this with configurable
decay curves:

\begin{equation}
\lambda(t) = \lambda_0 \cdot
\begin{cases}
  1 + (m_{\text{peak}} - 1) \cdot e^{-\alpha(T-t)} & \text{(exponential)} \\
  1 + (m_{\text{peak}} - 1) \cdot (1 - t/T)^{-\beta} & \text{(extended crunch)}
\end{cases}
\label{eq:period_end}
\end{equation}

where $T$ is the period-end date, $m_{\text{peak}}$ is the peak multiplier
(typically 3.5$\times$ for month-end, 6.0$\times$ for year-end), and
$\alpha, \beta$ are decay parameters.

\paragraph{Processing lags.}
Real posting dates lag event dates.  \system{} models this with log-normal
lag distributions per event type (e.g., goods receipt lag:
$\mu=1.5$~days, $\sigma=0.5$~days) and a cross-day posting model where the
probability of next-day posting increases with hour of origination.

\paragraph{Intraday patterns.}
Within each business day, transaction intensity follows configurable segments:
morning spike ($1.8\times$ at 08:30--10:00), lunch dip ($0.4\times$ at
12:00--13:30), and end-of-day rush ($1.5\times$ at 16:00--17:30).

\paragraph{Economic cycles and drift.}
Long-horizon datasets incorporate regime changes via:

\begin{equation}
\mu(t) = \mu_0 \cdot \left(1 + A\sin\!\left(\frac{2\pi t}{P}\right)\right)
  \cdot \mathbb{1}_{\text{expansion}}
  + \mu_0(1 - d_r) \cdot \mathbb{1}_{\text{recession}}
\label{eq:drift}
\end{equation}

where $P$ is the cycle period (default 48~months), $A$ is the amplitude
(0.15), and $d_r$ is the recession depth (0.25).  Regime switches are
modeled as Markov transitions with configurable probabilities.
